% !TEX root = task-6-1-6.tex
% Editable task bank. Append new BankTask blocks at the end of the document.
\documentclass[12pt,a4paper]{article}
\usepackage{physics-textbook}
\usepackage{graphicx,needspace}
\usetikzlibrary{decorations.pathmorphing}
\newif\ifShowSolutions
\ifdefined\TutorEdition
  \ShowSolutionstrue
\else
  \ShowSolutionsfalse % Change to \ShowSolutionstrue to display solutions.
\fi
\newcounter{banktask}
\newsavebox{\BankTaskBox}
\newcommand{\BankTaskContent}[3]{%
  \noindent\textbf{\thebanktask.}\quad #2\par
  \ifShowSolutions
    \Needspace{50mm}% Keep the solution heading with its opening text or diagram.
    #3\par
    \ifdefined\TutorEdition
      \ifcsname TutorSource#1\endcsname
        \nopagebreak\smallskip{\small\csname TutorSource#1\endcsname\par}%
      \fi
    \fi
  \fi
}
\newcommand{\BankTask}[3]{%
  \refstepcounter{banktask}\par\medskip
  \sbox{\BankTaskBox}{\begin{minipage}{\linewidth}%
    \BankTaskContent{#1}{#2}{#3}\end{minipage}}%
  \ifdim\dimexpr\ht\BankTaskBox+\dp\BankTaskBox\relax>\dimexpr\textheight-20mm\relax
    \Needspace{35mm}\BankTaskContent{#1}{#2}{#3}%
  \else
    \Needspace{\dimexpr\ht\BankTaskBox+\dp\BankTaskBox\relax}%
    \noindent\usebox{\BankTaskBox}\par
  \fi
  \medskip
}
\begin{document}
\BookSubsection{6.1.6}{Уравнение Эйнштейна для фотоэффекта}


% Source sdamgia-2301: https://phys-ege.sdamgia.ru/problem?id=2301
\BankTask{sdamgia-2301}{
Фотоэффект наблюдают, освещая поверхность металла светом фиксированной частоты. При этом задерживающая разность потенциалов равна U. После изменения частоты света задерживающая разность потенциалов увеличилась на $\Delta U=1{,}2\,\text{В}.$ На какую величину изменилась частота падающего света? (в $10^{14}$ Гц, округлив до десятых. Элементарный заряд  — $1{,}6\cdot10^{-19}$ Кл, постоянная Планка  — $6{,}6\cdot10^{-34}$ Дж·с.)
}{
\smallskip\noindent\textbf{Решение}\par
Запишем уравнение Эйнштейна для фотоэффекта для начальной частоты света

\[h\nu =A_{\,\text{вых}} + eU\]

и для измененной частоты

\[h ( \nu + \Delta \nu ) =A_{\,\text{вых}} + e ( U + \Delta U ) .\]

Вычтя из второго равенства первое, получим соотношение:

\[h\Delta \nu =e\Delta U \Longleftrightarrow \Delta \nu = \frac{e\Delta U}{h} = \frac{1{,}6 \cdot 10 ^{ - 19 } \,\text{Кл} \cdot 1{,}2\,\text{В}}{6{,}6 \cdot 10 ^{ - 34 } \,\text{Дж} \cdot \,\text{с}} \approx 2{,}9 \cdot 10 ^{ 14 } \,\text{Гц}.\]

Ответ: $2{,}9$.
}

% Source sdamgia-2302: https://phys-ege.sdamgia.ru/problem?id=2302
\BankTask{sdamgia-2302}{
Металлическую пластину освещают светом с энергией фотонов $6{,}2\,\text{эВ}$. Работа выхода для металла пластины равна $2{,}5\,\text{эВ}$. Какова максимальная кинетическая энергия образовавшихся фотоэлектронов? (в электрон-вольтах.)
}{
\smallskip\noindent\textbf{Решение}\par
Согласно уравнению Эйнштейна, энергия поглощенного фотона идет на работу выхода и на сообщение электрону кинетической энергии: $E_{\,\text{фот}}=A_{\,\text{вых}} + E_{\,\text{кин}}.$ Отсюда, для максимальной кинетической энергии образовавшихся фотоэлектронов имеем: 

\[E_{\,\text{кин}}=E_{\,\text{фот}} - A_{\,\text{вых}}=6{,}2\,\text{эВ} - 2{,}5\,\text{эВ}=3{,}7\,\text{эВ}.\]

Ответ: $3{,}7$.
}

% Source sdamgia-2303: https://phys-ege.sdamgia.ru/problem?id=2303
\BankTask{sdamgia-2303}{
Работа выхода электрона из металла $A_{\,\text{вых}}=3 \cdot 10 ^{ - 19 } \,\text{Дж}.$ Найдите максимальную длину волны $\lambda$ излучения, которым могут выбиваться электроны. (в нанометрах.) Постоянную Планка принять равной $6{,}6\cdot10^{-34}$ Дж·с, а скорость света  — $3\cdot10^{8}$ м/с.
}{
\smallskip\noindent\textbf{Решение}\par
Красная граница фотоэффекта определяется выражением $\nu _{\,\text{мин}}= \frac{A_{\,\text{вых}}}{h} .$ Принимая во внимание связь $\lambda \nu =c,$ для максимальной длины волны излучения, которым могут выбиваться электроны, имеем 

\[\lambda _{\max}= \frac{c}{\nu _{\,\text{мин}}} = \frac{ch}{A_{\,\text{вых}}} = \frac{3 \cdot 10 ^{ 8 } \,\text{м/с} \cdot 6{,}6 \cdot 10 ^{ - 34 } \,\text{Дж} \cdot \,\text{с}}{3 \cdot 10 ^{ - 19 } \,\text{Дж}} =660\,\text{нм}.\]

Ответ: $660$.
}

% Source sdamgia-2304: https://phys-ege.sdamgia.ru/problem?id=2304
\BankTask{sdamgia-2304}{
Поток фотонов с энергией $15\,\text{эВ}$ выбивает из металла фотоэлектроны, максимальная кинетическая энергия которых в $2$ раза меньше работы выхода. Какова максимальная кинетическая энергия образовавшихся фотоэлектронов? (в электрон-вольтах.)
}{
\smallskip\noindent\textbf{Решение}\par
Согласно уравнению Эйнштейна, энергия поглощенного фотона идет на работу выхода и на сообщение электрону кинетической энергии: $E_{\,\text{фот}}=A_{\,\text{вых}} + E_{\,\text{кин}}.$ По условию, $A_{\,\text{вых}}=2E_{\,\text{кин}}.$ Таким образом, для максимальной кинетической энергии фотоэлектронов имеем 

\[E_{\,\text{фот}}=3E_{\,\text{кин}} \Longleftrightarrow E_{\,\text{кин}}= \frac{1}{3} E_{\,\text{фот}}= \frac{1}{3} \cdot 15\,\text{эВ}=5\,\text{эВ}.\]

Ответ: $5$.
}

% Source sdamgia-2323: https://phys-ege.sdamgia.ru/problem?id=2323
\BankTask{sdamgia-2323}{
Работа выхода для материала катода вакуумного фотоэлемента равна $1{,}5\,\text{эВ}$. Катод освещается монохроматическим светом, у которого энергия фотонов равна $3{,}5\,\text{эВ}$. Чему равно запирающее напряжение, при котором фототок прекратится? (в вольтах.) Заряд электрона принять равным $1{,}6\cdot10^{-19}$ Кл, а $1\,\text{эВ}$  — $1{,}6\cdot10^{-19}$ Дж.
}{
\smallskip\noindent\textbf{Решение}\par
Запишем уравнение Эйнштейна для фотоэффекта $E_{\,\text{фот}}=A_{\,\text{вых}} + eU.$ Отсюда находим запирающее напряжение 

\[U= \frac{E_{\,\text{фот}} - A_{\,\text{вых}}}{e} = \frac{3{,}5\,\text{эВ} - 1{,}5\,\text{эВ}}{1{,}6 \cdot 10 ^{ - 19 } \,\text{Кл}} = \frac{2 \cdot 1{,}6 \cdot 10 ^{ - 19 } \,\text{Дж}}{1{,}6 \cdot 10 ^{ - 19 } \,\text{Кл}} =2\,\text{В}.\]

Ответ: $2$.
}

% Source sdamgia-3883: https://phys-ege.sdamgia.ru/problem?id=3883
\BankTask{sdamgia-3883}{
Энергия фотона, падающего на поверхность металлической пластинки, в $5$ раз больше работы выхода электрона с поверхности этого металла. Каково отношение максимальной кинетической энергии фотоэлектрона к работе выхода?
}{
\smallskip\noindent\textbf{Решение}\par
Согласно уравнению Эйнштейна, энергия фотона связана работой выхода и максимальной кинетической энергией соотношением: $h\nu=A_{\,\text{вых}} + E_{\,\text{кин}}.$ По условию, $h\nu=5A_{\,\text{вых}}.$ Следовательно, отношение максимальной кинетической энергии фотоэлектрона к работе выхода равно 

\[\frac{E_{\,\text{кин}}}{A_{\,\text{вых}}} = \frac{h\nu - A_{\,\text{вых}}}{A_{\,\text{вых}}} = \frac{5A_{\,\text{вых}} - A_{\,\text{вых}}}{A_{\,\text{вых}}} =4.\]

Ответ: $4$.
}

% Source sdamgia-3715: https://phys-ege.sdamgia.ru/problem?id=3715
\BankTask{sdamgia-3715}{
В таблице приведена зависимость максимальной кинетической энергии вылетающих из металла электронов от энергии падающих на металл фотонов.

\begin{center}\begin{tabular}{|c|c|c|c|c|}\hline
$E_{\,\text{фотона}},$ эВ&$2{,}4$&$2{,}8$&$3{,}3$&$4{,}0$\\\hline
$E_{\,\text{электрона}},$ эВ&$0{,}6$&$1{,}0$&$1{,}5$&$2{,}2$\\\hline
\end{tabular}\end{center}

Определите работу выхода для этого металла. (в электрон-вольтах.)
}{
\smallskip\noindent\textbf{Решение}\par
Согласно уравнению фотоэффекта Эйнштейна, энергия налетающего фотона идет на работу выхода и на сообщение электрону кинетической энергии: $E_{\,\text{фотона}}=A_{\,\text{выхода}} + E_{\,\text{электрона}}.$ Взяв данные любого столбца из таблицы, для работы выхода получаем:

\[A_{\,\text{выхода}}=E_{\,\text{фотона}} - E_{\,\text{электрона}}=2{,}4 \,\text{эВ} - 0{,}6\,\text{эВ}=1{,}8\,\text{эВ}.\]

Ответ: $1{,}8$.
}

% Source sdamgia-2321: https://phys-ege.sdamgia.ru/problem?id=2321
\BankTask{sdamgia-2321}{
В некоторых опытах по изучению фотоэффекта фотоэлектроны тормозятся электрическим полем. Напряжение, при котором поле останавливает и возвращает назад все фотоэлектроны, назвали задерживающим напряжением.
В таблице представлены результаты одного из первых таких опытов при освещении одной и той же пластины, в ходе которого было получено значение $h = 5{,}3 \cdot 10 ^{ - 34 } \,\text{Дж} \cdot \,\text{с}.$

\begin{center}\begin{tabular}{|c|c|c|}\hline
Задерживающее напряжение U, В&?&$0{,}6$\\\hline
Частота $\nu ,$ $10 ^{ 14 } \,\text{Гц}$&$5{,}5$&$6{,}1$\\\hline
\end{tabular}\end{center}

Чему равно опущенное в таблице первое значение задерживающего потенциала? (В вольтах; округлите до десятых.)
}{
\smallskip\noindent\textbf{Решение}\par
Согласно теории фотоэффекта, энергия поглощенного фотона идет на работу выхода и на сообщение электрону кинетической энергии. Электрическое поле совершает отрицательную работу, тормозя электроны. Таким образом, для первой частоты света и первого задерживающего напряжения имеем $h\nu _1=A_{\,\text{вых}} + eU_1,$ а для второй частоты и второго напряжения: $h\nu _2=A_{\,\text{вых}} + eU_2.$ Решая эту систему, принимая во внимание полученное в ходе эксперимента значение постоянной Планка, получаем выражение для первого значения задерживающего потенциала 

\[U_1=U_2 - \frac{h}{e} ( \nu _2 - \nu _1 ) =0{,}6\,\text{В} - \frac{5{,}3 \cdot 10 ^{ - 34 } \,\text{Дж} \cdot \,\text{с}}{1{,}6 \cdot 10 ^{ - 19 } \,\text{Кл}} ( 6{,}1 \cdot 10 ^{ 14 } \,\text{Гц} - 5{,}5 \cdot 10 ^{ 14 } \,\text{Гц} ) \approx 0{,}4\,\text{В}.\]

Ответ: $0{,}4\,\text{В}$.
}

% Source sdamgia-3005: https://phys-ege.sdamgia.ru/problem?id=3005
\BankTask{sdamgia-3005}{
Катод фотоэлемента с работой выхода $4{,}42 \cdot 10 ^{ - 19 } \,\text{Дж}$ освещается светом частотой $1{,}0 \cdot 10 ^{ 15 } \,\text{Гц}.$ Вылетевшие из катода электроны попадают в однородное магнитное поле с индукцией $8{,}3 \cdot 10 ^{ - 4 } \,\text{Тл}$ перпендикулярно линиям индукции этого поля. Чему равен максимальный радиус окружности R, по которой движутся электроны?
}{
\smallskip\noindent\textbf{Решение}\par
По уравнению Эйнштейна для фотоэффекта $h\nu = A + \frac{m v ^2}{2} .$
Электрон движется в магнитном поле под действием силы Лоренца $F=e v B.$ По второму закону Ньютона $F = ma.$ При этом электрон движется по окружности с центростремительным ускорением $a= \frac{v ^2}{R} .$
Объединяя формулы, выражаем радиус окружности, по которой движется электрон:

$R= \frac{\sqrt{2m ( h\nu - A )}}{eB} =$

\[R=\frac{\sqrt{2\cdot9\cdot10^{-31}\,\text{кг}\,(6{,}6\cdot10^{-34}\,\text{Дж}\cdot\text{с}\cdot10^{15}\,\text{Гц}-4{,}42\cdot10^{-19}\,\text{Дж})}}{1{,}6\cdot10^{-19}\,\text{Кл}\cdot8{,}3\cdot10^{-4}\,\text{Тл}}\approx5\cdot10^{-3}\,\text{м}=5\,\text{мм}.\]

Ответ: $5\,\text{мм}$.
}

% Source sdamgia-25424: https://phys-ege.sdamgia.ru/problem?id=25424
\BankTask{sdamgia-25424}{
Монохроматический свет попадает на литий, работа выхода которого равна $2{,}39\,\text{эВ}$ при длине волны $\lambda_0$ запирающее напряжение равно $U_3.$ Если изменить длину волны в $1{,}5$ раза, то необходимое запирающее напряжение увеличится вдвое. Найдите $\lambda_0$.
}{
\smallskip\noindent\textbf{Решение}\par
Запишем уравнение Эйнштейна для фотоэффекта $h v =A_{\,\text{вых}} + E_k.$ При этом частота света $v = \frac{c}{\lambda} ,$а максимальная кинетическая энергия электронов $E_k=eU.$
Зависимость задерживающего напряжения от длины волны обратная. Следовательно, при увеличении задерживающего напряжения в $2$ раза длина волны уменьшается в $1{,}5$ раза. Запишем уравнение фотоэффекта для обоих случаев:

$\frac{hc}{\lambda_0} =A_{\,\text{вых}} + eU_0;$

$\frac{1{,}5hc}{\lambda_0} =A_{\,\text{вых}} + 2eU_0.$

Решая эту систему уравнений, находим длину волны:

\[\lambda_0= \frac{hc}{2A_{\,\text{вых}}} = \frac{6{,}6 \cdot 10 ^{ - 34 } \cdot 3 \cdot 10 ^{8}}{2 \cdot 2{,}39 \cdot 1{,}6 \cdot 10 ^{ - 19 }} \approx 2{,}59 \cdot 10 ^{ - 7 } \,\text{м}=259\,\text{нм}.\]

Ответ: $259\,\text{нм}$.
}


% Source D4FFF9: https://ege.fipi.ru/bank/questions.php?search=1&proj=BA1F39653304A5B041B656915DC36B38&qid=D4FFF9
\BankTask{D4FFF9}{
Установите соответствие между графиками, представленными на рисунках, и зависимостями, которые они могут выражать.

\par\medskip\TaskMatchHeading{Графики}

А) рисунок 1

Б) рисунок 2

\par\medskip\TaskMatchHeading{Зависимости}

1) зависимость силы фототока от напряжения между электродами при неизменной освещённости

2) зависимость энергии фотона от длины волны

3) зависимость максимальной энергии фотоэлектронов от частоты света

4) зависимость энергии фотона от частоты света

\begin{tabular}{|c|c|}
\hline
А & Б \\
\hline
\phantom{2} & \phantom{3} \\
\hline
\end{tabular}

\medskip
\begin{center}
\begin{tikzpicture}[book diagram, x=0.024cm, y=0.024cm]
\draw[-{book tip}] (0,0) -- (150,0);
\draw[-{book tip}] (0,0) -- (0,127);
\node[below left] at (0,0) {$0$};
\draw[line width=0.5pt] (21,125) .. controls (22,118) and (21,105) .. (23,94) .. controls (26,73) and (32,58) .. (42,48) .. controls (54,36) and (69,29) .. (87,25) .. controls (104,21) and (125,19) .. (143,19);
\draw (0,103) -- (22,103);
\draw (22,103) -- (22,80);
\draw (22,32) -- (22,8);
\draw (0,23) -- (7,23);
\draw (32,23) -- (55,23);
\draw (80,23) -- (104,23);
\draw (104,23) -- (104,0);
\node[below] at (72,-18) {А) Рисунок 1};
\end{tikzpicture}
\hspace{1cm}
\begin{tikzpicture}[book diagram, x=0.024cm, y=0.024cm]
\draw[-{book tip}] (0,0) -- (163,0);
\draw[-{book tip}] (0,0) -- (0,127);
\node[below left] at (0,0) {$0$};
\draw[line width=0.7pt] (0,0) -- (52,0) -- (151,99);
\node[below] at (76,-18) {Б) Рисунок 2};
\end{tikzpicture}
\end{center}
}{
\smallskip\noindent\textbf{Решение от GPT}\par
График А — убывающая гипербола: $E=hc/\lambda$, поэтому ему соответствует зависимость 2. График Б описывает линейный рост максимальной энергии фотоэлектронов выше пороговой частоты: $K_{\max}=h\nu-A_{\text{вых}}$, то есть зависимость 3. Зависимость 1 представляет рост фототока с выходом на насыщение, а зависимость 4 — прямую $E=h\nu$, проходящую через начало координат; они не соответствуют данным графикам. Ответ: $23$.
}

% Source 369826: https://ege.fipi.ru/bank/questions.php?search=1&proj=BA1F39653304A5B041B656915DC36B38&qid=369826
\BankTask{369826}{
На металлическую пластинку падает пучок монохроматического света. При этом наблюдается явление фотоэффекта.

На графике А представлена зависимость энергии фотонов, падающих на катод, от физической величины \(x_1\), а на графике Б — зависимость максимальной кинетической энергии фотоэлектронов от физической величины \(x_2\).

Какая из физических величин отложена на горизонтальной оси на графике А и какая — на графике Б?

\TaskMatchingColumns[.53\linewidth]{Графики}{
А) график А\par
Б) график Б
}{Физическая величина \(x\)}{
1) частота\par
2) длина волны\par
3) массовое число\par
4) заряд ядра
}

\medskip
\begin{tabular}{|c|c|}
\hline
А & Б \\
\hline
\phantom{2} & \phantom{1} \\
\hline
\end{tabular}

\medskip
\begin{tikzpicture}[book diagram,x=0.45cm,y=0.23cm]
\node at (3.5,20.5) {А)};
\draw[-{book tip}] (0,0) -- (7.3,0) node[below] {$x_1$};
\draw[-{book tip}] (0,0) -- (0,18.5) node[left] {$E$};
\node[below left] at (0,0) {$0$};
\foreach \y in {0.5,1,...,17.5} {\draw (-0.09,\y) -- (0.09,\y);}
\foreach \y in {5,10,15} {\node[left] at (-0.14,\y) {$\y$};}
\foreach \xx/\yy in {1/15,2/7.5,3/5,4/3.75,5/3,6/2.5} {
\draw[thin] (0,\yy) -- (\xx,\yy) -- (\xx,0);
}
\draw[thick] plot[domain=0.82:7,samples=100] (\x,{15/\x});
\foreach \x/\y in {1/15,2/7.5,3/5,4/3.75,5/3,6/2.5}{\fill (\x,\y) circle[radius=1pt];}
\end{tikzpicture}
\qquad
\begin{tikzpicture}[book diagram,x=0.45cm,y=0.45cm]
\node at (3.5,5.25) {Б)};
\draw[-{book tip}] (0,0) -- (7.1,0) node[below] {$x_2$};
\draw[-{book tip}] (0,-1.25) -- (0,4.7) node[left] {$E$};
\node[left] at (0,0) {$0$};
\draw[thick] (0,0) -- (2,0) -- (5.7,3.7);
\end{tikzpicture}
}{
\smallskip\noindent\textbf{Решение от GPT}\par
Энергия фотона \(E=h\nu=hc/\lambda\), поэтому убывающая гипербола А соответствует длине волны (2). По уравнению Эйнштейна \(K_{\max}=h\nu-A_{\text{вых}}\), поэтому возрастающая прямая Б соответствует частоте (1). Массовое число (3) и заряд ядра (4) не являются переменными в этих зависимостях. Ответ: А — 2, Б — 1.
}

\end{document}
